% Part: first-order-logic
% Chapter: syntax-and-semantics
% Section: covered-structures

\documentclass[../../../include/open-logic-section]{subfiles}

\begin{document}

\olfileid{fol}{syn}{cov}

\olsection{\printtoken{P}{structure} المغطاة للغات الرتبة الأولى}

\begin{explain}
قد تقدم أن الحد الذي يخلو من ال!!{variable}s حد \emph{مغلق}.
\end{explain}

\begin{defn}[!!^{value} الحدود المغلقة]
إذا كان $t$ حدًا مغلقًا من اللغة~$\Lang L$، وكانت $\Struct M$
!!a{structure} لـ~$\Lang L$، عُرّفت \emph{!!{value}}~$\Value{t}{M}$
كما يأتي:
\begin{enumerate}
\item إذا كان $t$ هو ال!!{constant}~$c$ فحسب، فإن
  $\Value{c}{M} = \Assign{c}{M}$.
\item إذا كان $t$ من الصورة $\Atom{f}{t_1, \ldots, t_n}$، فإن
  \[
  \Value{t}{M} = \Assign{f}{M}(\Value{t_1}{M}, \ldots,
  \Value{t_n}{M}).
  \]
\end{enumerate}
\end{defn}

\begin{defn}[ال!!{structure} المغطاة]
نسمي ال!!{structure} \emph{مغطاة} متى كان لكل عنصر من المجال
حد مغلق يكون ذلك العنصر !!{value} له.
\end{defn}

\begin{ex}
لنأخذ~$\Lang L$ لغةً ذات ال!!{constant}s~$\Obj{zero}$،
و$\Obj{one}$، و$\Obj{two}$، و\dots، وفيها ال!!{predicate} الثنائية~$<$،
وال!!{function}s الثنائية~$+$ و$\times$. ثم نعيّن~$\Struct M$
!!a{structure} لهذه اللغة~$\Lang L$، مجالها
$\Domain M = \{0, 1, 2, \ldots \}$، والتعيينات فيها
$\Assign{\Obj{zero}}{M} = 0$، و$\Assign{\Obj{one}}{M} = 1$،
و$\Assign{\Obj{two}}{M} = 2$، وعلى هذا القياس. أما رمز العلاقة
الثنائية~$<$، فنأخذ~$\Assign{<}{M}$ مجموعة الأزواج
$\tuple{c_1, c_2} \in \Domain{M}^2$ التي يكون فيها $c_1$ أصغر
من~$c_2$. فمن أمثلتها $\tuple{1, 3} \in \Assign{<}{M}$، وأما
$\tuple{2, 2} \notin \Assign{<}{M}$. وأما
ال!!{function} الثنائية~$+$، فنجعل $\Assign{+}{M}$ الجمع
المعتاد؛ ومنه $\Assign{+}{M}(2,3)=5$. وعلى هذا نؤول
ال!!{function} الثنائية~$\times$ بالضرب المعتاد. فتكون !!{value}
~$\Obj{four}$ هي~$4$، وتُحسب !!{value}
~$\times(\Obj{two}, +(\Obj{three},\Obj{zero}))$، وهي بالترميز
الوسطي $\Obj{two} \times (\Obj{three} + \Obj{zero})$، كما يأتي:
\begin{multline*}
\Value{\times(\Obj{two}, +(\Obj{three},\Obj{zero}))}{M} =\\
\begin{aligned}
& \Assign{\times}{M}(\Value{\Obj{two}}{M}, \Value{+(\Obj{three}, \Obj{zero})}{M})\\
& = \Assign{\times}{M}(\Value{\Obj{two}}{M}, \Assign{+}{M}(\Value{\Obj{three}}{M},
\Value{\Obj{zero}}{M})) \\
& = \Assign{\times}{M}(\Assign{\Obj{two}}{M}, \Assign{+}{M}(\Assign{\Obj{three}}{M},
\Assign{\Obj{zero}}{M})) \\
& = \Assign{\times}{M}(2, \Assign{+}{M}(3, 0)) \\
& = \Assign{\times}{M}(2, 3) \\
& = 6
\end{aligned}
\end{multline*}
\end{ex}

\begin{prob}
هل $\Struct N$، وهو النموذج القياسي للحساب، بنية مغطاة؟ بيّن وجه جوابك.
\end{prob}

\end{document}
